\section{Fault Detection and Protection Logic}
\label{sec:ch3-protection-logic}

The protection manager receives calibrated voltage, calibrated RMS current,
estimated socket temperature, and TinyML arc output. It then assigns warning
and trip states according to thresholds and persistence rules in the firmware.
The relay is opened only after a trip condition is accepted by the protection
logic; warnings alone do not disconnect the load.

\subsection{Voltage Supervision}
\label{sec:ch3-voltage-protection}

Voltage supervision uses calibrated RMS voltage and mains-present state. The
firmware defines the normal operating band as 200-250V. Instantaneous
undervoltage protection is set at 170V, delayed undervoltage is set below
200V for 7s, delayed overvoltage is set above 250V for 7s, and
instantaneous overvoltage is set at 260V. Recovery requires the voltage to
return to 212-250V and remain stable for 5min.

These thresholds are treated as design choices for supply-state supervision
rather than validated voltage-protection performance, since the experimental
validation focused on fire-inducing fault types. The implemented values are
summarized in \textit{Table~\ref{tab:ch3-voltage-thresholds}}.

\begin{table}[H]
  \centering
  \caption{Voltage-protection thresholds verified from firmware}
  \label{tab:ch3-voltage-thresholds}
  \begin{tabular}{@{}p{0.31\textwidth}p{0.17\textwidth}p{0.44\textwidth}@{}}
    \toprule
    \textbf{Condition} & \textbf{Firmware value} & \textbf{Interpretation} \\
    \midrule
    Normal lower limit & 200V & Lower edge of normal operation \\
    Normal upper limit & 250V & Upper edge of normal operation \\
    Instant undervoltage & 170V & Immediate severe undervoltage trip \\
    Delayed undervoltage & 200V for 7s & Persistent low-voltage trip \\
    Delayed overvoltage & 250V for 7s & Persistent high-voltage trip \\
    Instant overvoltage & 260V & Immediate severe overvoltage trip \\
    Recovery band & 212-250V for 5min & Stable voltage required before recovery \\
    \bottomrule
  \end{tabular}
\end{table}

\textit{Table~\ref{tab:ch3-voltage-thresholds}} shows that the voltage logic includes
both instantaneous and delayed decisions. This avoids treating every brief
voltage excursion as a latched trip while still allowing severe voltage
conditions to disconnect promptly.

\subsection{Overload Protection}
\label{sec:ch3-overload-protection}

The overload logic separates warning from relay trip. A warning begins when
current exceeds 10.0A after the warning delay and clears below 9.6A. A
sustained-overload trip requires the current average to remain above 14.5A
long enough for the overload score to reach the trip score.

The persistence score is represented as
\begin{equation}
  S_{\mathrm{OL}}[k] =
  \min\left(S_{\max}, S_{\mathrm{OL}}[k-1] + 1\right),
  \qquad I_{\mathrm{avg}}\geq 14.5\mathrm{A},
  \label{eq:overload-score-rise}
\end{equation}
and
\begin{equation}
  S_{\mathrm{OL}}[k] =
  \max\left(0, S_{\mathrm{OL}}[k-1] - 1\right),
  \qquad I_{\mathrm{avg}}\leq 14.1\mathrm{A}.
  \label{eq:overload-score-fall}
\end{equation}
In (\ref{eq:overload-score-rise}) and (\ref{eq:overload-score-fall}),
$S_{\mathrm{OL}}$ is the sustained-overload score, $I_{\mathrm{avg}}$ is the
firmware's 1s averaged current, and $S_{\max}=6$. The score is updated every
10s, and relay trip occurs at $S_{\mathrm{OL}}=6$. Therefore, a persistent
overload near or above 14.5A trips after about 60s, while shorter excursions
can produce warnings without disconnection. Persistence-based behavior is
consistent with the principle that overcurrent protection should distinguish
brief transients from sustained thermal stress [22].

\subsection{Socket-Heating Protection}
\label{sec:ch3-heating-protection}

Socket-heating protection uses the estimated socket temperature
$T_{\mathrm{socket,est}}$ from (\ref{eq:estimated-socket-temp}). The firmware
uses a warning threshold of $39.0\degreeC$, a warning-clear threshold of
$38.6\degreeC$, and a trip threshold of $40.0\degreeC$. The warning hysteresis
reduces rapid warning toggling near the threshold, while the trip threshold
creates a conservative response to possible contact heating.

The trip rule is
\begin{equation}
  H =
  \begin{cases}
  1, & T_{\mathrm{socket,est}}\geq 40.0\degreeC,\\
  0, & T_{\mathrm{socket,est}}<40.0\degreeC .
  \end{cases}
  \label{eq:heating-trip-rule}
\end{equation}
In (\ref{eq:heating-trip-rule}), $H$ is the heating-trip state. This trip
criterion is intentionally based on the compensated socket estimate rather
than the direct device temperature measurement because the external thermistor can under-read the
contact hotspot during degraded-contact heating [30], [46], [47].

\subsection{Series Arc-Fault Protection}
\label{sec:ch3-arc-protection}

The Random Forest arc model produces a frame-level arc class. The protection
manager does not trip from a single isolated model output. Instead, the
firmware uses a leaky counter:
\begin{equation}
  C_{\mathrm{arc}}[k] =
  \min\left(16, C_{\mathrm{arc}}[k-1] + 2\right),
  \qquad \hat{y}_{\mathrm{arc}}=1,
  \label{eq:arc-counter-rise}
\end{equation}
and
\begin{equation}
  C_{\mathrm{arc}}[k] =
  \max\left(0, C_{\mathrm{arc}}[k-1] - 6\right),
  \qquad \hat{y}_{\mathrm{arc}}=0 .
  \label{eq:arc-counter-fall}
\end{equation}
In (\ref{eq:arc-counter-rise}) and (\ref{eq:arc-counter-fall}),
$C_{\mathrm{arc}}$ is the arc persistence counter and
$\hat{y}_{\mathrm{arc}}$ is the TinyML model output. The relay trips when
$C_{\mathrm{arc}}\geq 8$. The counter therefore requires repeated arc-positive
frames and quickly decays when the classifier output returns to normal.

The firmware also includes a current-return verification path. If current
falls below 0.03A before the arc counter reaches the trip threshold, the
system waits up to 200ms to verify whether current returns before treating the
event as an arc trip. This guards against disconnecting on a simple load turn
off while still allowing fast response to arc-like interruption and restrike.

\subsection{Protection Priority and Latching}
\label{sec:ch3-priority-latching}

When multiple protection conditions are present, the firmware resolves the
reported fault according to protection priority. Socket heating is given the
highest priority because temperature is the most direct indicator of thermal
fire risk, followed by arc fault, sustained overload, and voltage-related
protection states. The relay-latching circuit is pulsed by the firmware after a
trip decision, and the fault state remains latched until the user clears it.
This prioritization does not change the underlying detection equations; it only
determines which fault reason is reported when conditions overlap.

The principal thresholds are summarized in \textit{Table~\ref{tab:ch3-protection-summary}}.
\begin{table}[H]
  \centering
  \caption{Principal protection thresholds verified from firmware}
  \label{tab:ch3-protection-summary}
  \begin{tabular}{p{0.26\textwidth}p{0.30\textwidth}p{0.34\textwidth}}
    \toprule
    \textbf{Protection function} & \textbf{Warning or model threshold} & \textbf{Relay-trip condition} \\
    \midrule
    Voltage supervision & Outside 200--250V band & Instant 170/260V or delayed 7s low/high voltage \\
    Overload & Warning above 10.0A; clear below 9.6A & 1s current average at or above 14.5A for about 60s \\
    Socket heating & Warning at 39.0°C; clear at 38.6°C & Estimated socket temperature at or above 40.0°C \\
    Series arc fault & RF probability threshold depends on context & Arc counter reaches 8 from repeated arc-positive frames \\
    \bottomrule
  \end{tabular}
\end{table}

\textit{Table~\ref{tab:ch3-protection-summary}} shows that the prototype uses a hybrid
method: deterministic thresholds for voltage, overload, and temperature, and
TinyML-supported classification for series arcs. The hybrid design is intended
to complement existing protection devices by observing socket-level thermal
and waveform behavior that may not be sufficient to trip a branch breaker.

